\subsection{A power-contraction refinement of the published mixed rounding}

This is an unconditional written corollary of the construction and certified
estimates in Heilman--Jones--Malavolta,
\href{https://arxiv.org/html/2609.25462v1}{Sections 5--6}.
The new polynomial, suffix and arithmetic statements are also proved in
\path{Grothe/Haagerup/ScalarTailRefinement.lean}. The published Gaussian
interpretation and integration facts are explicit inputs to the Lean score
theorem; the complete Grothendieck conclusion is not yet formalized here.
The fixed phase-tail estimate used below is independently derived in Lean,
so it is not an extra hypothesis of the final score corollary.

Let $\rho_5(t)=tP_5(t^2)$ be exactly the degree-75 preprocessor in their
equation (5.4), with positive raw selector weight
$w_5=0.0000169358532612748746$. Its 38 rational coefficients are stored in
\path{data/scalar-tail-refinement/certificate.json}. Exact convolution gives
\[
 \|P_5\|_1=\frac{999999999993}{10^{12}}\le1,
 \qquad \|P_5^{16}\|_1<\frac{581}{1000}=:q.
\]
The latter has degree 592 and is checked by one outward-rounded native
calculation, independently replayed with integer arithmetic.

For every $j$, quotient and remainder give
\[
 \|P_5^{2j+1}\|_1\le
 q^{\lfloor(2j+1)/16\rfloor}.
\]
In particular, for $j\ge67$ the exponent is at least 8. Multiplication by
the monomial in $\rho_5^{2j+1}$ does not change its Wiener norm. Using the
published Wallis estimate $\sum_{j\ge67}h_j\le1/268$, the entire omitted
outer part of branch 5 costs at most
\[
 \frac{w_5q^8}{268}
 =8.20505217878206\ldots\,10^{-10}.
\]
For the formal fixed phase-tail estimate, let $r_j=4h_j/\pi$ and
$g_j=r_j(j+1)(2j+1)$. The existing telescoping majorant and an exact
rational check give
\[
 \sum_{j\ge67}h_j\le\frac{\pi}{4}
 \left(\sum_{j=67}^{119}r_j+\frac{g_{120}}{240}\right)
 \le\frac{355}{452}
 \left(\sum_{j=67}^{119}r_j+\frac{g_{120}}{240}\right)
 \le\frac1{268}.
\]
The finite comparison is kernel-checked, with no additional native fact.
This replaces only equation (6.20)'s coarse reserve $w_5/268$. All output
degrees of the first 67 outer blocks remain paid, including their complete
tail beyond output index 119. The omitted-outer bound may also pay degrees
inside the retained prefix; this only makes the estimate conservative.

The published lower slope and upper costs, equations (6.2), (6.19) and
(6.30), give the exact rational sufficient budget
\begin{align*}
 a&=0.71179640805640433452,\\
 R&=1.074829404498064\,10^{-6}
   +2.699176355797438\,10^{-12}
   +1.040224125770398\,10^{-7}
   +\frac{w_5q^8}{268},\\
 M&=\frac{1250000000000000121997}{1250000000000000000000}.
\end{align*}
Here $M$ is the exact selector mass from (6.28); it is not rounded to one.
Rational arithmetic proves
\[
 \frac{a-R}{M}>\frac1{1.40489844},
 \qquad \frac{M}{a-R}<1.404898431637415.
\]
The valid-kernel residual-transfer argument of the cited paper therefore
yields
\[
 \boxed{K_G^{\mathbb C}<1.40489844.}
\]

All four pinned non-regeneration verifiers were executed successfully with
python-flint 0.9.0. They check polynomial nonvanishing, directed endpoint-
subtracted energy integration, the coefficient head and every tail budget,
and final normalization. The two accepted coefficient tables were checked
against the published source hashes; they were not regenerated in this run.
The input provenance, replay logs and exact arithmetic are recorded in
\path{data/scalar-tail-refinement/PROOF.md} and its verification directory.
There is no new comparison or Hardy conjecture in this written deduction.
The repository's separate fully formal unconditional bound remains
$K_G^{\mathbb C}\le1.404909543$ until the published analytic inputs are
connected in Lean.
